% \iffalse
%<*driver>
\ProvidesFile{seqenseigne-core-eval.dtx}%
[2026/05/08 v1.0-dev seqenseigne – Module core-eval (évaluations sommatives)]
%</driver>
% \fi
%
% =============================================================
% seqenseigne-core-eval.dtx
% Module : CORE-EVAL
% Rôle   : macros et environnements liés aux évaluations sommatives :
%          - page de titre d'évaluation (\seqTitreEval)
%          - barème (en-tête de l'évaluation, environnement seqEvalBareme)
%          - exercices d'évaluation (environnement seqEvalExercice)
%          - corrigés différés (\seqEvalCorrigeExo, \seqEvalAfficheCorriges)
%
% Module créé en v0.13.5.2 par extraction depuis seqenseigne-core.dtx.
%
% v0.13.5.2 étape 2 : renommages pour cohérence du préfixe seq.
%   - \titreEval               → \seqTitreEval
%   - \corrigeExoEval          → \seqEvalCorrigeExo
%   - \seqEvalQCMItem          → \seqQcmQuestion (déplacée vers core-exos)
%   - compteur ExoEvalNum      → seqExoEvalNum (centralisé dans core)
%
% v0.13.5.2 étape 3 : refonte fonctionnelle.
%   \seqTitreEval :
%   - Suppression des options numSeq, titreSeq, niveauDiff (et leur
%     mapping vers un nombre de points "Ce sujet permet d'obtenir N
%     points"). Une éval peut couvrir plusieurs séquences, ou un seul
%     objectif d'une séquence — pas d'attache à une séquence unique.
%   - Ajout de l'option titre (texte libre) en remplacement de
%     numSeq+titreSeq.
%   - Suppression du bandeau supérieur "Séquence NN" et de l'indication
%     "Niveau fondamental/avancé/exploration" à droite.
%   - "Mise en pratique" → "Évaluation" dans le rendu.
%   - Conservation de theme=, etab=, classe=.
%
%   seqEvalExercice :
%   - Option "nom" renommée en "titre" (cohérence avec \seqTitreEval).
%   - Ajout de l'option "bareme" : si renseignée, affichage en italique
%     aligné à droite du cartouche, format [/ <bareme>] (ex. [/ 5 points]).
%
% v0.13.5.2 étape 4 : nouvel environnement seqEvalObjectifs.
%   Tableau d'objectifs évalués avec grille de critères de maîtrise
%   (Insuffisant / À consolider / Satisfaisant / Très bon, à cocher
%   manuellement par l'enseignant). Affiché par l'appli uniquement si
%   le mode de notation inclut les critères ("criteres" ou
%   "note_criteres"). Items écrits en clair par l'utilisateur :
%     <code> & <intitulé> & \ding{113} & \ding{113} & \ding{113} & \ding{113} \\
%   (tabularray refuse les macros qui contiennent des & et \\).
%
% Le QCM, étant un format particulier d'exercice utilisable hors-éval
% (en livret de séquence par exemple), n'est pas défini ici. Il est
% défini dans core-exos comme environnement enfant utilisable dans
% n'importe quel contexte (seqEvalExercice, seqExercice).
%
% Dépendances :
%  - seqenseigne-core (boîtes \begin{boiteTitreGen}, \begin{boiteTitreSection},
%    \begin{boitePaleVariable}, \seq@setColors, \eue, \@charlb, \@charrb,
%    \newwrite\corrfileeval, compteur seqExoEvalNum)
%  - seqenseigne-core-exos (environnement seqQcm pour les QCM en éval)
%  - paquet pifont (déjà chargé via dinglist pour le QCM) pour les
%    cases à cocher \ding{113} dans seqEvalObjectifs
%
% =============================================================
%
%<*package>
\NeedsTeXFormat{LaTeX2e}%
\ProvidesPackage{seqenseigne-core-eval}%
[2026/05/08 v1.0-dev seqenseigne core-eval module]

%% ------------------------------------------------------------------
%% ---- Traitement des options --------------------------------------
%% ------------------------------------------------------------------
\newif\ifdocstd\docstdtrue
\DeclareOption{ltxdoc}{\docstdfalse}%
\ProcessOptions\relax

%% Note : tous les compteurs (seqExoEvalNum, seqQcmItemNum, etc.) sont
%% déclarés dans seqenseigne-core (centralisation des compteurs depuis
%% v0.13.5.2 étape 2).

%% ------------------------------------------------------------------
%% ---- \seqInitCorrigesEval ----------------------------------------
%% ------------------------------------------------------------------
%% v0.13.5.6 — Allocation du canal d'écriture des corrigés d'éval.
%% Pendant exact de \seqInitCorriges (core-exos, canal \corrfile) et
%% \seqInitAnnexes (core, canal \annfile). Encapsule l'allocation
%% \newwrite\corrfileeval qui était auparavant au top-level de
%% seqenseigne-core.sty — les allocations TeX top-level ne sont pas
%% indexées par le parser de paquet côté appli (services/
%% paquet_parseur.py), et la fermeture transitive du préambule
%% atomique les perdait. La macro publique, elle, est correctement
%% indexée comme \newcommand.
%%
%% À appeler avant \seqTitreEval. Côté appli, services/
%% render_evaluation.py::_section_titre l'émet systématiquement.
\newcommand{\seqInitCorrigesEval}{%
\newcounter{seqExoEvalNum}%
\newwrite\corrfileeval
}%

%% ------------------------------------------------------------------
%% ---- \seqTitreEval -----------------------------------------------
%% ------------------------------------------------------------------
%    \begin{macro}{\seqTitreEval}%
%% v0.13.5.2 étape 3 — Refonte fonctionnelle :
%%   - Suppression de numSeq, titreSeq, niveauDiff (et leur mapping
%%     vers un nombre de points). Une éval peut couvrir plusieurs
%%     séquences ou un seul objectif d'une séquence — le titre est
%%     ce que l'enseignant veut.
%%   - Ajout de "titre" (texte libre) à la place de numSeq+titreSeq.
%%   - Suppression du bandeau supérieur "Séquence NN" (plus de
%%     référence à une séquence unique).
%%   - Suppression de l'indication "Niveau fondamental/avancé/..."
%%     à droite (n'a plus de sens pour une éval qui mélange
%%     potentiellement plusieurs séries).
%%   - Suppression du bloc "Ce sujet permet d'obtenir N points" qui
%%     dépendait de niveauDiff.
%%   - "Mise en pratique" → "Évaluation".
%%   - Conservation de theme=, etab=, classe=.
%%   - Le sous-titre du corrigé devient simplement "Corrigé".
%%
%% Note : on garde les noms de groupes d'options (titreEval, exerciceEval)
%% en interne, sans préfixe seq, conformément à la convention validée :
%% préfixe seq pour macros et environnements publics, pas pour les
%% groupes d'options xkeyval.
\define@cmdkey[seq]{titreEval}{titre}{}%
\define@cmdkey[seq]{titreEval}{etab}{}%
\define@cmdkey[seq]{titreEval}{classe}{}%
\define@cmdkey[seq]{titreEval}{theme}[gris]{}%
%% ---- Page de titre de l'évaluation ------------------------------
\newcommand{\seqTitreEval}[1][]{%
	%% ---- Options
	\setkeys[seq]{titreEval}{theme,#1}%
	%% ---- Ouverture du fichier des corrigés (le \newwrite est déclaré
	%%      globalement dans seqenseigne-core)
	\immediate\openout\corrfileeval=CorrList.tex
	%% ---- Début du corrigé : titre
	%% Note : \begin{boiteTitreGen} prend un argument obligatoire (le
	%% bandeau supérieur). Comme on ne veut plus de bandeau "Séquence NN"
	%% en étape 3, on lui passe un argument vide via \unexpanded{{}}.
	%% C'est l'équivalent côté write de \begin{boiteTitreGen}{} ligne 130.
	\immediate\write\corrfileeval{%
		\noexpand\cleardoublepage
		\noexpand\begin{boiteTitreGen}\unexpanded{{}}%
			\unexpanded{\begin{minipage}[c][70pt][c]{\linewidth}}%
				\unexpanded{\vfill \textit{\'Evaluation: corrig\'e}}%
				\noexpand\vfill \noexpand\centering \@charlb
				\noexpand\Huge \noexpand\textbf{\eue{\cmdseq@titreEval@titre}} \@charrb
				\noexpand\vfill \noexpand\hfill
				\noexpand\textit{\eue{\cmdseq@titreEval@etab} - \eue{\cmdseq@titreEval@classe}}%
				\unexpanded{ \end{minipage}	\end{boiteTitreGen}}%
}%
%% ---- Entêtes et pieds de page
\ifdocstd
	\pagestyle{fancy}%
	\fancyhf{}%
	\lfoot{Le \today}%
	\rfoot{Page \thepage/\pageref{LastPage}}%
\fi
%% ---- Couleurs
\seq@setColors{\cmdseq@titreEval@theme}%
%% ---- Boîte de titre
\begin{boiteTitreGen}{}%
	\begin{minipage}[c][70pt][c]{\linewidth}%
		\vfill \textit{Évaluation}%
		\vfill \centering {\Huge \textbf{\cmdseq@titreEval@titre}}%
		\vfill	\hfill
		\textit{\cmdseq@titreEval@etab~- \cmdseq@titreEval@classe}%
	\end{minipage}%
\end{boiteTitreGen}%
\begin{center} \begin{boitePaleVariable}{\linewidth}%
		Toutes les réponses doivent être justifiées,
		sauf si une indication contraire est donnée. \newline
		Pour chaque question, si le travail n'est pas terminé,
		laisser tout de même une trace de la recherche. \newline
		Elle sera prise en compte dans la notation.
	\end{boitePaleVariable} \end{center}%
}%
%    \end{macro}%
% -------------------------------------------------------------

%% ---- Pour chaque corrigé -----------------------------------
%    \begin{macro}{\seqEvalCorrigeExo}%
%% Ecrit le corrigé de l'exercice dans un fichier séparé
\newcommand{\seqEvalCorrigeExo}[1]{%
	\scantokens{%
		\begin{tcbverbatimwrite}{Corriges/c\theseqExoEvalNum.tex}%
			#1
		\end{tcbverbatimwrite}%
	}%
}%
%    \end{macro}%
% -------------------------------------------------------------

%% ---- Regroupe les corrigés en fin d'évaluation -------------
%    \begin{macro}{\seqEvalAfficheCorriges}%
%% qui commence sur une page impaire
\newcommand{\seqEvalAfficheCorriges}{%
	\immediate\closeout\corrfileeval
	\clearpage
	\ifodd\c@page\else
		\null\newpage
		\if@twocolumn\null\newpage\fi
	\fi

	 \clearpage \seqTitreSectionCorrige { Corrigés des exercices de la série exploration } \label {corrigeExos3} 
 \par \nopagebreak \vspace {5pt}\nopagebreak \begin {boiteTitreExo} \par Exercice\;2\par \par \textit {Des caractères sacrés}\end {boiteTitreExo} \par \nopagebreak \vspace {5pt}\nopagebreak   \begin {seqColEnum}[nbCols=1] \item Je sais que $1~302 = 1\times 1~000 + 3 \times 100 + 2 \times 1$. Or, dans le premier nombre en hiéroglyphes, je vois : \begin {seqColItem}[nbCols=1] \item une seule fois le premier hiéroglyphe, \includegraphics [scale=0.2]{N09S01E02_corrige4.png}, donc il vaut 1~000. \item trois fois le deuxième hiéroglyphe, \includegraphics [scale=0.2]{N09S01E02_corrige3.png}, donc il vaut 100. \item deux fois le troisième hiéroglyphe, \includegraphics [scale=0.2]{N09S01E02_corrige1.png}, donc il vaut 1. \end {seqColItem} Je sais que $10~047 = 1 \times 10~000 + 4 \times 10 + 7 \times 1$. Or, dans le deuxième nombre en hiéroglyphes, je vois : \begin {seqColItem}[nbCols=1] \item une seule fois le premier hiéroglyphe, \includegraphics [scale=0.2]{N09S01E02_corrige5.png}, donc il vaut 10~000. \item quatre fois le deuxième hiéroglyphe, \includegraphics [scale=0.2]{N09S01E02_corrige2.png}, donc il vaut 10. \item sept fois le le troisième hiéroglyphe, \includegraphics [scale=0.2]{N09S01E02_corrige1.png} que je connais déjà. \end {seqColItem} Je sais que $2~311~021 = 2 \times 1~000~000 + 3 \times 100~000 + 1 \times 10~000 + 1 \times 1~000 + 2 \times 10 + 1 \times 1$. \newline Or, dans le troisième nombre en hiéroglyphes, je vois : \begin {seqColItem}[nbCols=1] \item deux fois le premier hiéroglyphe, \includegraphics [scale=0.2]{N09S01E02_corrige7.png}, donc il vaut 1~000~000. \item trois fois le deuxième hiéroglyphe, \includegraphics [scale=0.2]{N09S01E02_corrige6.png}, donc il vaut 100~000. \item une seule fois le troisième hiéroglyphe, \includegraphics [scale=0.2]{N09S01E02_corrige5.png} que je connais déjà. \item une seule fois le quatrième hiéroglyphe, \includegraphics [scale=0.2]{N09S01E02_corrige4.png} que je connais déjà. \item deux fois le cinquième hiéroglyphe, \includegraphics [scale=0.2]{N09S01E02_corrige2.png} que je connais déjà. \item une seule fois le sixième hiéroglyphe, \includegraphics [scale=0.2]{N09S01E02_corrige1.png} que je connais déjà. \end {seqColItem} Voici donc dans l'ordre les valeurs des sept hiéroglyphes : \begin {seqColItem}[nbCols=3] \item \includegraphics [scale=0.2]{N09S01E02_corrige1.png} vaut 1. \newline \item \includegraphics [scale=0.2]{N09S01E02_corrige2.png} vaut 10. \newline \item \includegraphics [scale=0.2]{N09S01E02_corrige3.png} vaut 100. \newline \item \includegraphics [scale=0.2]{N09S01E02_corrige4.png} vaut 1~000. \newline \item \includegraphics [scale=0.2]{N09S01E02_corrige5.png} vaut 10~000. \newline \item \includegraphics [scale=0.2]{N09S01E02_corrige6.png} vaut 100~000. \newline \item \includegraphics [scale=0.2]{N09S01E02_corrige7.png} vaut 1~000~000. \end {seqColItem} \item En m'aidant de la question précédente, je peux calculer les nombres pour chaque hiéroglyphe : \newline \includegraphics [scale=0.2]{N09S01E02_enonce4.png} \quad vaut $3 \times 10 + 2 \times 1 = 32$ \newline \includegraphics [scale=0.2]{N09S01E02_enonce5.png} \quad vaut $1 \times 100~000 + 3 \times 10~000 + 2 \times 1 = 130~002$ \newline \includegraphics [scale=0.2]{N09S01E02_enonce6.png} \quad vaut $1 \times 1~000~000 + 2 \times 1~000 + 1 \times 10 = 1~002~010$ \item Je peux écrire $426 = 4 \times 100 + 2 \times 10 + 6 \times 1$ et donc \includegraphics [scale=0.9]{N09S01E02_corrige8.png}. \newline Je peux écrire $527 = 5 \times 100 + 2 \times 10 + 7 \times 1$ et donc \includegraphics [scale=0.9]{N09S01E02_corrige9.png}. \newline Je peux écrire $12~315 = 1 \times 10~000 + 2 \times 1~000 + 3 \times 100 + 1 \times 10 + 5 \times 1$ et donc \includegraphics [scale=0.9]{N09S01E02_corrigeA.png}. \newline Je peux écrire $1~234~000 = 1 \times 1~000~000 + 2 \times 100~000 + 3 \times 10~000 + 4 \times 1~000$ et donc \includegraphics [scale=0.9]{N09S01E02_corrigeB.png}. \newline \end {seqColEnum}  
 
%
}%
%    \end{macro}%
% -------------------------------------------------------------

%% ------------------------------------------------------------------
%% ---- Environnement seqEvalBareme ---------------------------------
%% ------------------------------------------------------------------
%    \begin{environment}{seqEvalBareme}%
%% v0.13.5.2 étape 4 patch v2 :
%%   La macro \seqEvalBaremeItem a été supprimée pour la même raison
%%   que \seqEvalObjectifsItem et \seqQcmItem : tabularray refuse
%%   les macros qui contiennent des & (séparateurs de cellules).
%%   L'utilisateur (ou l'appli génératrice) écrit chaque ligne en
%%   clair :
%%
%%     \begin{seqEvalBareme}%
%%         Exercice I & 4 points \\
%%         Exercice II & 6 points \\
%%         Présentation et usage de la langue française & 1 point \\
%%     \end{seqEvalBareme}%
%%
%%   colspec={X[5,l,m] X[r,m]} : la colonne libellé est large et
%%   alignée à gauche, la colonne points est étroite et alignée à
%%   droite. Plus besoin de \hfill explicite : tabularray gère
%%   l'alignement via la spec de colonne.
\newenvironment{seqEvalBareme}{%
	\begin{center} \begin{boitePaleVariable}{0.8\linewidth}%
			\begin{tblr}{%
					colspec = {X[5,l,m] X[r,m]},
					rowsep = 2pt,
				}%
				\SetCell[c=2]{c} \large \textbf{Barème} & \\
				}{%
			\end{tblr}%
		\end{boitePaleVariable} \end{center}%
}%
%    \end{environment}%
% -------------------------------------------------------------

%% ------------------------------------------------------------------
%% ---- Environnement seqEvalObjectifs ------------------------------
%% ------------------------------------------------------------------
%    \begin{environment}{seqEvalObjectifs}%
%% v0.13.5.2 étape 4 — Nouveauté :
%%   Tableau d'objectifs évalués avec grille de critères de maîtrise.
%%   Affiché uniquement si l'enseignant a choisi un mode de notation
%%   qui inclut les critères ("criteres" ou "note_criteres" côté appli).
%%   Pour les modes "note" ou "aucun", l'appli ne génère pas ce
%%   tableau — il n'a pas de raison d'être sans critères de maîtrise.
%%
%%   Les 4 colonnes de cases à cocher correspondent aux niveaux de
%%   maîtrise atteints : Insuffisant, À consolider, Satisfaisant,
%%   Très bon. L'enseignant coche manuellement après correction.
%%
%% Limitation tabularray (cf. seqQcm pour le même genre de contrainte) :
%%   tabularray refuse les macros qui contiennent des & ou des \\ dans
%%   leur corps. Donc PAS DE MACRO \seqEvalObjectifsItem qui
%%   produirait toute la ligne. L'utilisateur (ou l'appli génératrice)
%%   écrit chaque ligne en clair :
%%
%%     \begin{seqEvalObjectifs}%
%%         Obj. 02 & Pour les nombres décimaux\ldots
%%             & \ding{113} & \ding{113} & \ding{113} & \ding{113} \\
%%         Obj. 04 & Calculer le PGCD\ldots
%%             & \ding{113} & \ding{113} & \ding{113} & \ding{113} \\
%%     \end{seqEvalObjectifs}%
%%
%%   La verbosité des 4 \ding{113} répétés est sans importance en
%%   pratique : c'est l'appli qui génère ce code, pas l'enseignant
%%   à la main.
%%
%%   Le code dans la première colonne est un agrégat fourni par
%%   l'appli, par exemple "Obj. 02" si tous les objectifs viennent
%%   de la même séquence, ou "Séq. S03 Obj. 13" pour identifier
%%   la séquence d'origine.
%%
%%   colspec figé en X[2,c,m] X[5,l,m] X[1,c,m] X[1,c,m] X[1,c,m] X[1,c,m]
%%   (limitation tabularray : pas de macro à l'intérieur).

\newenvironment{seqEvalObjectifs}{%
	\begin{center} \begin{boitePaleVariable}{\linewidth}%
			\begin{tblr}{%
					colspec={X[2,c,m] X[5,l,m] X[1,c,m] X[1,c,m] X[1,c,m] X[1,c,m]},
					rowsep = 2pt,
				}%
				\SetCell[c=6]{c} \large \textbf{Niveaux de maîtrise des objectifs évalués} & & & & & \\
				\toprule
				\SetCell[c=2]{c} \textbf{Objectif} &
					& \rotatebox{90}{\textit{Insuffisant}}%
					& \rotatebox{90}{\textit{À consolider}}%
					& \rotatebox{90}{\textit{Satisfaisant}}%
					& \rotatebox{90}{\textit{Très bon}} \\
				\midrule
				}{%
			\end{tblr}%
		\end{boitePaleVariable} \end{center}%
}%
%    \end{environment}%
% -------------------------------------------------------------

%% ------------------------------------------------------------------
%% ---- Environnement seqEvalExercice -------------------------------
%% ------------------------------------------------------------------
%    \begin{environment}{seqEvalExercice}%
%% v0.13.5.2 étape 3 — Refonte fonctionnelle :
%%   - Option "nom" renommée en "titre" (cohérence avec \seqTitreEval).
%%   - Ajout d'une option "bareme" (texte libre type "5 points",
%%     "3,5 points"). Si renseignée, affichage en italique aligné à
%%     droite du cartouche, format [/ <bareme>].
%% ---- Options
\define@cmdkey[seq]{exerciceEval}{titre}{}%
\define@cmdkey[seq]{exerciceEval}{bareme}{}%
\newenvironment{seqEvalExercice}[1][]{%
	%% ---- Réinitialisation des options pour permettre la détection
	%%      de leur absence éventuelle dans cet appel
	\let\cmdseq@exerciceEval@titre\@undefined
	\let\cmdseq@exerciceEval@bareme\@undefined
	%% ---- Options
	\setkeys[seq]{exerciceEval}{#1}%
	%% ---- Numéro de l'exercice
	\stepcounter{seqExoEvalNum}%
	%% ---- Début du corrigé
	\immediate\write\corrfileeval{%
		\unexpanded{\par\nopagebreak\vspace{5pt}\nopagebreak}%
		\unexpanded{\begin{boiteTitreSection}}%
				\noexpand\begin{minipage}[c][20pt][c]{\noexpand\linewidth}%
					\noexpand\Large \noexpand\textbf{ Exercice \Roman{seqExoEvalNum}%
						\ifx\cmdseq@exerciceEval@titre\@undefined\else : \eue{\cmdseq@exerciceEval@titre} \fi }%
					\unexpanded{\end{minipage} \end{boiteTitreSection}}%
		\unexpanded{\par\nopagebreak\vspace{5pt}\nopagebreak}%
		\noexpand\input{Corriges/c\theseqExoEvalNum.tex}%
	}%
	%% ---- Boite de titre
	\par\vspace{10pt}%
	\begin{boiteTitreSection}%
		\begin{minipage}[c][20pt][c]{\linewidth}%
			\Large \textbf{Exercice \Roman{seqExoEvalNum}%
				\ifx\cmdseq@exerciceEval@titre\@undefined\else : \cmdseq@exerciceEval@titre\fi }%
			\ifx\cmdseq@exerciceEval@bareme\@undefined\else
				\hfill \textit{[/ \cmdseq@exerciceEval@bareme]}%
			\fi
		\end{minipage}%
	\end{boiteTitreSection}%
	\par\vspace{10pt}%
}{}%
%    \end{environment}%
% -------------------------------------------------------------

\endinput
%</package>
%
% \Finale
\endinput
